\documentclass[11pt]{article}
% Shared preamble for the rigor documents (Lamport-style structured proofs).  \input this file after \documentclass.
\usepackage[T1]{fontenc}\usepackage{lmodern}\usepackage{microtype}
\usepackage[margin=1in]{geometry}
\usepackage{amsmath,amssymb,amsthm,mathtools}
\usepackage{xcolor}\usepackage{enumitem}\usepackage{booktabs}\usepackage{graphicx}\usepackage{hyperref}
\usepackage[most]{tcolorbox}
\definecolor{navy}{HTML}{17365D}\definecolor{teal}{HTML}{117A8B}\definecolor{warn}{HTML}{A33A2B}\definecolor{okgreen}{HTML}{2E7D4F}
\hypersetup{colorlinks=true,linkcolor=navy,citecolor=teal,urlcolor=teal}
\theoremstyle{plain}
\newtheorem{theorem}{Theorem}[section]\newtheorem{lemma}[theorem]{Lemma}\newtheorem{proposition}[theorem]{Proposition}\newtheorem{corollary}[theorem]{Corollary}
\theoremstyle{definition}\newtheorem{definition}[theorem]{Definition}\newtheorem{remark}[theorem]{Remark}\newtheorem{claim}[theorem]{Claim}\newtheorem{issue}[theorem]{Issue}
% ---------- Lamport structured proofs ----------
% \lstep{level}{number}{assertion}   -> indented "<level>number. assertion"
% \lproof                             -> "PROOF:" line (start of the sub-proof of the preceding step)
% \lby{justification}                 -> "BY ..." leaf justification for the preceding step
% \lqed{level}{number}                -> QED step of a sub-proof
% keywords: \lassume \lprove \lsuffices \lcase \ldefine \lpick \lnew
\newlength{\lindent}\setlength{\lindent}{1.7em}
\newcommand{\lstep}[3]{\par\smallskip\noindent\hspace*{\dimexpr\numexpr#1-1\relax\lindent\relax}{\color{navy}$\langle#1\rangle#2$.}\ #3}
\newcommand{\lproof}{\par\noindent\hspace*{\lindent}{\scshape Proof:}}
\newcommand{\lby}[1]{\par\noindent\hspace*{\lindent}{\scshape By}\ #1.}
\newcommand{\lqed}[2]{\par\smallskip\noindent\hspace*{\dimexpr\numexpr#1-1\relax\lindent\relax}{\color{navy}$\langle#1\rangle#2$.}\ {\scshape Q.E.D.}}
\newcommand{\lassume}{{\scshape Assume}\ }\newcommand{\lprove}{{\scshape Prove}\ }\newcommand{\lsuffices}{{\scshape Suffices}\ }
\newcommand{\lcase}{{\scshape Case}\ }\newcommand{\ldefine}{{\scshape Define}\ }\newcommand{\lpick}{{\scshape Pick}\ }\newcommand{\lnew}{{\scshape New}\ }
% ---------- verdict boxes ----------
\newtcolorbox{verdictok}{colback=okgreen!8,colframe=okgreen,title={\bfseries Verdict: verified},fonttitle=\small}
\newtcolorbox{verdictgap}{colback=orange!10,colframe=orange!80!black,title={\bfseries Verdict: gap / caveat},fonttitle=\small}
\newtcolorbox{verdictbreak}{colback=warn!10,colframe=warn,title={\bfseries Verdict: proof-breaking issue},fonttitle=\small}
\newtcolorbox{citedbox}[1]{colback=navy!5,colframe=navy!60,title={\bfseries Cited result: #1},fonttitle=\small,breakable}
\setlength{\parindent}{0pt}\setlength{\parskip}{4pt}\allowdisplaybreaks
\newcommand{\cA}{\mathcal A}\newcommand{\cF}{\mathcal F}\newcommand{\cM}{\mathfrak M}\newcommand{\cN}{\mathfrak N}

% extra calligraphic shortcuts (\providecommand: no clash if a document defines its own)
\providecommand{\cB}{\mathcal B}\providecommand{\cC}{\mathcal C}\providecommand{\cD}{\mathcal D}\providecommand{\cE}{\mathcal E}\providecommand{\cG}{\mathcal G}\providecommand{\cH}{\mathcal H}\providecommand{\cI}{\mathcal I}\providecommand{\cJ}{\mathcal J}\providecommand{\cK}{\mathcal K}\providecommand{\cL}{\mathcal L}\providecommand{\cO}{\mathcal O}\providecommand{\cP}{\mathcal P}\providecommand{\cQ}{\mathcal Q}\providecommand{\cR}{\mathcal R}\providecommand{\cS}{\mathcal S}\providecommand{\cT}{\mathcal T}\providecommand{\cU}{\mathcal U}\providecommand{\cV}{\mathcal V}\providecommand{\cW}{\mathcal W}\providecommand{\cX}{\mathcal X}\providecommand{\cY}{\mathcal Y}\providecommand{\cZ}{\mathcal Z}
\newcommand{\Fid}{\operatorname{F}}\newcommand{\Tr}{\operatorname{Tr}}\newcommand{\eps}{\varepsilon}\newcommand{\dd}{\,\mathrm d}

\newcommand{\hh}{\mathfrak h}
\DeclareMathOperator{\Var}{Var}
\DeclareMathOperator{\diag}{diag}
\DeclareMathOperator*{\argmax}{arg\,max}
\newcommand{\e}[1]{\times10^{#1}}
\DeclareMathOperator{\Li}{Li}
\DeclareMathOperator{\sgn}{sgn}
\DeclareMathOperator{\sech}{sech}
\newcommand{\Zmax}{Z_{\max}}
\DeclareMathOperator{\spec}{spec}
\newcommand{\Erec}{E_{\mathrm{rec}}}

\title{SDP certificate numerics for the quasi-free recovery programme\\
{\large Track S11 (numerical), Phase 3b}}
\author{Claude (track S11)}
\date{2026-09-10}

\begin{document}
\maketitle

\begin{abstract}
Numerical verification of the KKT/duality structure of the second-order quasi-free
recovery programme
$\min_{(X,Z)\in F} g_Q(Q-Q_{\rm rec}(X,Z))$
of \texttt{optimality\_second\_order.tex} (S8) and \texttt{optimality\_all\_channels.tex}.
Part (a): on the lattice chain $(L_A,L_B,L_C)=(4,8,2)$ we compute the symmetric
logarithmic derivative (SLD) $t_{\rm opt}$ of S8's optimal defect, maximise the
\emph{linear} functional $\Tr(t_{\rm opt}Q_{\rm rec}(X,Z))$ over the feasible set $F$
by an SDP, and test the KKT identity and strong duality
$\Delta(t_{\rm opt})^2/(8V(t_{\rm opt}))=g_Q(\delta_{\rm opt})$.
Part (b): in the continuum modular-frame discretisation of the compression defect we
compute the SLD $t_c$, evaluate $\Delta(t_c)$ by an SDP over $F$, and report the
dual ratio $\Delta(t_c)^2/(8V(t_c))\big/(f_2\zeta^2)$, $f_2=c/(12\pi^2)$.
\end{abstract}


\section{Conventions}
\label{sec:conv}

\paragraph{The programme.}  Throughout, $\hh_I=\hh_A\oplus\hh_D$, $\hh_D=\hh_B\oplus\hh_C$,
$Q$ is the vacuum one-particle symbol on $\hh_I$ ($Q_{jk}=\omega(a^*(e_k)a(e_j))$, $0<Q<1$),
and, following S8 (\texttt{optimality\_second\_order.tex} Thm.~3.1) and
\texttt{optimality\_all\_channels.tex} Thm.~3.5,
\begin{equation}\label{eq:F}
 \mathcal F:=\Bigl\{(X,Z):\ X:\hh_B\to\hh_D,\ Z=Z^*\ \text{on}\ \hh_D,\
   XQ_{BB}X^*\le Z\le 1-X(1-Q_{BB})X^*\Bigr\},\quad
 Q^{\rm rec}(X,Z)=\begin{pmatrix}Q_{AA}&Q_{AB}X^*\\ XQ_{BA}&Z\end{pmatrix}.
\end{equation}
Both constraints are entered into the solver as Schur complements, exactly as in
\texttt{relaxation\_test.py}: with $Q_{BB}=\sigma\sigma^*$ and $1-Q_{BB}=\tau\tau^*$
(Cholesky),
\begin{equation}\label{eq:schur}
 \begin{pmatrix}Z&X\sigma\\ \sigma^*X^*&1_{\hh_B}\end{pmatrix}\ge0,
 \qquad
 \begin{pmatrix}1_{\hh_D}-Z&X\tau\\ \tau^*X^*&1_{\hh_B}\end{pmatrix}\ge0 .
\end{equation}

\paragraph{The objective and its Legendre dual.}  Write $Q=U\diag(q)U^*$ and
$\widetilde\delta:=U^*\delta U$.  For a weight matrix $w=(w_{ik})$, $w_{ik}=w_{ki}>0$, put
\begin{equation}\label{eq:gw}
 g_w(\delta):=\tfrac14\sum_{ik}w_{ik}\,|\widetilde\delta_{ik}|^2,
 \qquad
 V_w(t):=\sum_{ik}\frac{|\widetilde t_{ik}|^2}{2\,w_{ik}} .
\end{equation}
The \emph{true} quantum-Fisher weights are $w^{\rm true}_{ik}=\bigl(q_i(1-q_k)+q_k(1-q_i)\bigr)^{-1}$;
S8 also uses the \emph{capped} weights $w^{(\kappa_c)}:=\min\{w^{\rm true},1+\cosh\kappa_c\}$
(S8 Sec.~3.4), which preserve convexity and still lower-bound the true form.

\begin{lemma}[Legendre pair]\label{lem:legendre}
For every self-adjoint $\delta$ and every $w$ as above,
$g_w(\delta)=\tfrac18\sup_{t=t^*}(\Tr t\delta)^2/V_w(t)$, the supremum being attained at the
\emph{$w$-SLD} $t=\tfrac12U(w\circ\widetilde\delta)U^*$ ($\circ$ = entrywise product), for which
\begin{equation}\label{eq:sldid}
 \Tr t\,\delta=2\,g_w(\delta),\qquad V_w(t)=\tfrac12 g_w(\delta),\qquad
 \frac{(\Tr t\delta)^2}{8V_w(t)}=g_w(\delta).
\end{equation}
Moreover $V_{w^{\rm true}}(t)=\Tr\bigl(tQt(1-Q)\bigr)=\Var_\omega\dd\Gamma(t)$, but for a
\emph{capped} $w$ one has $V_w(t)\ne\Tr(tQt(1-Q))$ in general (indeed $V_w\ge V_{w^{\rm true}}$).
\end{lemma}

\begin{proof}
\lstep{1}{1}{$\Tr t\delta=\sum_{ik}\widetilde t_{ik}\overline{\widetilde\delta_{ik}}$.}
\lby{$\Tr t\delta=\Tr\widetilde t\widetilde\delta=\sum_{ik}\widetilde t_{ik}\widetilde\delta_{ki}$ and
$\widetilde\delta_{ki}=\overline{\widetilde\delta_{ik}}$ ($\delta=\delta^*$)}
\lstep{1}{2}{$|\Tr t\delta|^2\le\bigl(\sum_{ik}|\widetilde t_{ik}|^2/w_{ik}\bigr)\bigl(\sum_{ik}w_{ik}|\widetilde\delta_{ik}|^2\bigr)
 =2V_w(t)\cdot4g_w(\delta)=8V_w(t)g_w(\delta)$.}
\lby{$\langle1\rangle1$ and Cauchy--Schwarz with the weights $w_{ik}^{\pm1/2}$}
\lstep{1}{3}{Equality holds iff $\widetilde t_{ik}=\lambda\,w_{ik}\widetilde\delta_{ik}$ for a
constant $\lambda>0$; $\lambda=\tfrac12$ gives \eqref{eq:sldid} by direct substitution.}
\lby{the equality case of Cauchy--Schwarz; substitution:
$\Tr t\delta=\tfrac12\sum w|\widetilde\delta|^2=2g_w$ and
$V_w(t)=\sum\tfrac14w^2|\widetilde\delta|^2/(2w)=\tfrac18\sum w|\widetilde\delta|^2=\tfrac12g_w$}
\lstep{1}{4}{$\Tr(tQt(1-Q))=\sum_{ik}|\widetilde t_{ik}|^2q_k(1-q_i)
 =\sum_{ik}|\widetilde t_{ik}|^2\tfrac12\bigl(q_k(1-q_i)+q_i(1-q_k)\bigr)=V_{w^{\rm true}}(t)$.}
\lby{expand in the $Q$-eigenbasis; symmetrise in $(i,k)$ using $|\widetilde t_{ik}|=|\widetilde t_{ki}|$}
\lqed{1}{5}
\end{proof}

\noindent Consequently the dual value of \texttt{optimality\_all\_channels.tex} eq.~(4.1),
\begin{equation}\label{eq:dual}
 \Delta(t):=\inf_{(X,Z)\in\mathcal F}\Tr\,t\bigl(Q-Q^{\rm rec}(X,Z)\bigr),\qquad
 \mathcal D:=\sup_t\frac{\Delta(t)_+^2}{8V(t)},\qquad V=V_{w^{\rm true}},
\end{equation}
satisfies $\mathcal D=\min_{\mathcal F}g_{w^{\rm true}}$ (Thm.~4.1(3), Sion), and the same
identity holds verbatim with $(g_w,V_w)$ in place of $(g_{w^{\rm true}},V)$ for any capped $w$.
\emph{The pairing must be consistent}: mixing $g_{w^{(\kappa_c)}}$ with $V_{w^{\rm true}}$ produces
a spurious duality gap, and this is visible in Table~\ref{tab:lattice} below.



\section{Part (a): lattice KKT at S8's optimum}
\label{sec:lattice}

\paragraph{Set-up.}  Half-filled hopping chain, $(L_A,L_B,L_C)=(4,8,2)$, $n=14$,
$Q_{jk}=\sin(\pi(j-k)/2)/(\pi(j-k))$, $z=L_AL_C/(L_B(L_A+L_B+L_C))=1/14$,
$2f_2z^2=8.615747\times10^{-5}$ (two chiralities, $c=1$; the reference value of S8 Table~1).
The spectrum of $Q$ obeys $\min\{q,1-q\}=5.711\times10^{-10}$, i.e.\
$\kappa_{\max}=\log\frac{1-q_{\min}}{q_{\min}}=21.283$ and $w^{\rm true}_{\max}=8.755\times10^{8}$:
eight orders of magnitude of weight, which is why S8 caps.  Script:
\texttt{numerics/optimality\_all/kkt\_lattice.py}, output \texttt{kkt\_lattice.out}
(cvxpy 1.7 + Clarabel, SCS fallback at $\eps=10^{-11}$; complex variables $X\in M_{10\times8}$,
$Z=Z^*\in M_{10}$, two $18\times18$ LMIs \eqref{eq:schur}).

\paragraph{What is computed, for each cap $\kappa_c$ of the ladder.}
(i) the primal $g_{w}^{\min}:=\min_{\mathcal F}g_w$ by SDP;
(ii) the $w$-SLD $t_{\rm opt}$ of the minimising defect $\delta_{\rm opt}$;
(iii) $\Delta(t_{\rm opt})=\inf_{\mathcal F}\Tr t_{\rm opt}(Q-Q^{\rm rec})$ by a second (linear) SDP;
(iv) the two dual ratios $\Delta^2/(8V_w)/g^{\min}_w$ and $\Delta^2/(8V_{\rm true})/g^{\min}_w$;
(v) the Frobenius distance between the recovered symbol of the \emph{linear} maximiser and that
of the primal minimiser (the KKT statement proper);
(vi) for reference, $g_w$ evaluated at S8's stored $(X,Z)$ of \texttt{best\_4\_8\_2.npz}.

\begin{table}[h]\centering\small
\begin{tabular}{rlllllll}
\toprule
$\kappa_c$&$\min_{\mathcal F}g_w$&$g_w$ at S8's $(X,Z)$&$\Delta(t_{\rm opt})$&
$\dfrac{2g_w^{\min}-\Delta}{2g_w^{\min}}$&$\dfrac{\Delta^2/(8V_w)}{g_w^{\min}}$&
$\dfrac{\Delta^2/(8V_{\rm true})}{g_w^{\min}}$&$\|\delta_{\rm lin}-\delta_{\rm opt}\|_F$\\
\midrule
 4    &$6.38797\e{-5}$&$6.72726\e{-5}$&$1.27759\e{-4}$&$+6.3\e{-6}$&$0.9999873$&$1.052656$&$5.7\e{-4}$\\
 6    &$7.32364\e{-5}$&$7.37944\e{-5}$&$1.46472\e{-4}$&$+7.7\e{-6}$&$0.9999846$&$1.012533$&$2.1\e{-3}$\\
 8    &$7.69024\e{-5}$&$7.72397\e{-5}$&$1.53804\e{-4}$&$+3.2\e{-6}$&$0.9999936$&$1.023669$&$1.4\e{-3}$\\
10    &$8.21078\e{-5}$&$8.42693\e{-5}$&$1.64222\e{-4}$&$-3.7\e{-5}$&$1.0000747$&$1.034610$&$3.8\e{-3}$\\
12    &$9.48589\e{-5}$&$1.13953\e{-4}$&$1.89658\e{-4}$&$+3.2\e{-4}$&$0.9993682$&$1.073697$&$2.2\e{-2}$\\
14    &$1.10426\e{-4}$&$2.21125\e{-4}$&$2.20830\e{-4}$&$+9.7\e{-5}$&$0.9998062$&$1.061769$&$2.0\e{-2}$\\
16    &$1.26865\e{-4}$&$9.51051\e{-4}$&$2.54204\e{-4}$&$-1.9\e{-3}$&$1.0037407$&$1.044722$&$1.2\e{-2}$\\
18    &$1.35239\e{-4}$&$2.86823\e{-3}$&$2.70801\e{-4}$&$-1.2\e{-3}$&$1.0023866$&$1.018609$&$5.6\e{-2}$\\
21.283&$1.19686\e{-4}$&$4.99549\e{-2}$&$2.59240\e{-4}$&$-8.3\e{-2}$&$1.1728910$&$1.1728910$&$3.2\e{-1}$\\
\bottomrule
\end{tabular}
\caption{Lattice $(4,8,2)$.  Cap ladder $w^{(\kappa_c)}=\min(w^{\rm true},1+\cosh\kappa_c)$; the last
row $\kappa_c=21.283=\kappa_{\max}$ is the \emph{uncapped} problem.  Column~6 is the strong-duality
ratio with the \emph{consistent} Legendre pair \eqref{eq:gw}; column~7 is the same with
$V_{\rm true}=\Tr(tQt(1-Q))$, which is the correct pairing only in the last row (where columns 6
and 7 coincide to $10^{-8}$, as Lemma~\ref{lem:legendre} demands).  Rows $\kappa_c\ge12$ are
reported by Clarabel as \texttt{optimal\_inaccurate}; the last two rows are solver-limited
(a ratio $>1$ is mathematically impossible).}
\label{tab:lattice}
\end{table}


\paragraph{Findings (a).}
\begin{enumerate}[leftmargin=1.6em,itemsep=2pt]
\item \emph{Strong duality holds numerically.}  With the consistent Legendre pair, column~6 of
Table~\ref{tab:lattice} is $1$ to $1.5\times10^{-5}$ for $\kappa_c\le10$ and to
$4\times10^{-3}$ up to $\kappa_c=16$; the deviation tracks the solver status, not the
mathematics.  This is a direct numerical confirmation of
\texttt{optimality\_all\_channels.tex} Thm.~4.1(3) (the Sion minimax step
$\min_{\mathcal K}\sup_t=\sup_t\min_{\mathcal K}$) in finite dimensions.
\item \emph{KKT holds at the optimum.}  The relative KKT gap
$\bigl(2g_w^{\min}-\Delta(t_{\rm opt})\bigr)/(2g_w^{\min})$ (column~5) is
$3\times10^{-6}$--$3\times10^{-4}$ for $\kappa_c\le14$: the SLD of the optimal defect, used as a
\emph{linear} functional, is minimised over $\mathcal F$ at the optimum itself.  The last column
shows that the linear programme's minimiser reproduces the primal recovered symbol to
$10^{-3}$--$10^{-2}$ in Frobenius norm.  The residual is expected and is \emph{not} a failure of
KKT: the linear programme's minimiser need not be unique (the objective is flat along the
face of $\mathcal F$ on which the optimum sits --- at $(4,8,2)$ eight of the ten singular values
of $X$ are pinned at $1$, so both LMIs are active on an $8$-dimensional face), and only the value
$\Delta(t_{\rm opt})$, not the argument, is certified by duality.
\item \emph{The cap must be paired with its own variance.}  Column~7 exceeds $1$ by
$1$--$7\%$ throughout the capped ladder.  This is not a violation of Theorem~4.1: capping changes
the quadratic form, and the quantity $\Delta(t)^2/(8\Tr(tQt(1-Q)))$ is a valid lower bound for
$\min_{\mathcal F}g_{w^{\rm true}}$, not for $\min_{\mathcal F}g_{w^{(\kappa_c)}}$.  The
inequality $V_{w^{(\kappa_c)}}\ge V_{w^{\rm true}}$ of Lemma~\ref{lem:legendre} accounts for the
sign.  In the uncapped row columns 6 and 7 agree to $2\times10^{-8}$, as they must.
\item \emph{S8's stored $(X,Z)$ is a good but not exact minimiser of the capped form.}
At its own cap $\kappa_c=8$ it gives $g_w=7.72397\times10^{-5}$ against the SDP optimum
$7.69024\times10^{-5}$: $0.44\%$ high, consistent with the log-barrier tolerance of
\texttt{opt\_gaussian.py}.  (Its \emph{exact} error $-\log\Fid=8.6695\times10^{-5}$, S8 Table~1,
is unaffected: that number is a rigorous upper bound evaluated at the point actually produced.)
Away from $\kappa_c=8$ the stored point degrades fast --- at $\kappa_c=21.28$ it gives
$g_{w^{\rm true}}=5.0\times10^{-2}$, i.e.\ $420\times$ the true optimum: the capped optimiser
does dump defect into the ultraviolet, exactly as S8's Sec.~3.4 predicts.
\item \emph{The uncapped lattice value is a lattice artefact, as S8 says.}  Reading down
column~2, $\min_{\mathcal F}g_{w^{\rm true}}$ has not converged at $\kappa_c=18$
($1.35\times10^{-4}$, i.e.\ $1.57\times2f_2z^2$) and the solver fails beyond.  At $n=14$ the
modular cutoff $\kappa_{\max}=21.3$ is not large compared with $\log(1/z)=2.6$, so the ``order of
limits'' obstruction of S8 Sec.~3.4 is active; the continuum programme of Section~\ref{sec:continuum}
is the correct place to compare with $f_2\zeta^2$.
\end{enumerate}

\begin{verdictok}[\;(a) Strong duality and KKT are numerically verified on the lattice]
For the convex programme $\min_{\mathcal F}g_w$ on $(4,8,2)$, and for every cap
$\kappa_c\in\{4,6,8,10,12,14\}$ at which Clarabel converges, the dual value
$\Delta(t_{\rm opt})^2/(8V_w(t_{\rm opt}))$ equals the primal optimum to
$\le3\times10^{-4}$ relative, and the linear programme $\inf_{\mathcal F}\Tr t_{\rm opt}(Q-Q^{\rm rec})$
attains its value at (a point with the same value as) the primal optimum.  The SLD of the optimal
defect is therefore the dual certificate, which is the structure that
\texttt{optimality\_all\_channels.tex} Thm.~4.1(3) asserts and that track S10 must exhibit
analytically at the compression.  No inconsistency with Thm.~4.1 was found; the two apparent
violations (column~7, and the $\kappa_c\ge16$ rows of column~6) are respectively an inconsistent
Legendre pairing and interior-point solver failure, both explained above.
\end{verdictok}


\section{Part (b): continuum compression defect}
\label{sec:continuum}

\subsection{A block-exact modular-frame discretisation}
\label{sec:disc}
Geometry: $A=(-a,0)$, $B=(0,b)$, $C=(b,L)$, $D=BC=(0,L)$, $b=L/(1+s)$,
$\zeta=as/(a+L)$ (equal to the cross ratio $z$ of the four points), $a=1$, $L=2$.
Modular coordinate of the interval $AD=(-a,L)$:
\begin{equation}\label{eq:xi}
 \xi=\log\frac{x+a}{L-x},\qquad x(\xi)=\frac{Le^\xi-a}{1+e^\xi},\qquad
 \frac{dx}{d\xi}=\beta(x)=\frac{(x+a)(L-x)}{a+L},
\end{equation}
so that $A=\{\xi<\xi_0\}$, $B=\{\xi_0<\xi<\xi_1\}$, $C=\{\xi>\xi_1\}$ with
$\xi_0=\log(a/L)$, $\xi_1=\log\frac{a+b}{L-b}$.  In half-density form the vacuum symbol is the
translation-invariant kernel of Note~5 / \texttt{defect\_matrix.py},
\begin{equation}\label{eq:qt}
 \widetilde q(t)=\tfrac12\delta(t)-\frac{i}{4\pi\sinh(t/2)},\qquad t=\xi-\xi',
 \qquad\text{Fourier symbol } \frac{1}{1+e^{\kappa}},\ \kappa=2\pi k ,
\end{equation}
and the zero-collar compression defect $\delta_c=Q-Q^{\rm rec}(X_c,Z_c)$ is supported on the
$A\!-\!D$ blocks (Moebius invariance of the vacuum forces $Z_c=Q_{DD}$, i.e.\ $\delta_{c,DD}=0$),
\begin{equation}\label{eq:dc}
 \widetilde\delta_c(\xi,\xi')=\frac{i}{2\pi}\sqrt{\beta(x)\beta(y)}\;R_s(-x,y)+\text{h.c.},\qquad
 R_s(u,v)=\frac{s\,u\,v}{(u+v)\bigl(L(u+v)+s\,u\,v\bigr)},
\end{equation}
$x=x(\xi)<0<y=x(\xi')$ (kernel \texttt{exact}); the $O(s)$ part is $R_s^{(1)}=\frac{s}{L}\frac{uv}{(u+v)^2}$
(kernel \texttt{first}).

\begin{lemma}[Closed form of the discretised vacuum symbol]\label{lem:closed}
Let $\{I_m=[\alpha_m,\beta_m]\}$ be disjoint intervals of the $\xi$-line, $h_m=|I_m|$, and
$e_m=h_m^{-1/2}\chi_{I_m}$.  Put
\(
 F_2(t):=-4\bigl(\Li_2(-e^{-|t|})-\Li_2(e^{-|t|})+\tfrac{\pi^2}{4}\bigr)\sgn t .
\)
Then $F_2$ is odd, $F_2''(t)=1/\sinh(t/2)$, $F_2(0)=0$, $F_2(\pm\infty)=\mp\pi^2$, and
\begin{equation}\label{eq:Qcell}
 \langle e_m,Qe_n\rangle=\tfrac12\delta_{mn}
 -\frac{i}{4\pi\sqrt{h_mh_n}}\Bigl[F_2(\beta_m-\alpha_n)-F_2(\alpha_m-\alpha_n)
 -F_2(\beta_m-\beta_n)+F_2(\alpha_m-\beta_n)\Bigr].
\end{equation}
\end{lemma}

\begin{proof}
\lstep{1}{1}{$\int_0^X\log\tanh u\,du=-\tfrac12\bigl(\Li_2(-e^{-2X})-\Li_2(e^{-2X})+\tfrac{\pi^2}{4}\bigr)$
for $X>0$.}
\lby{$\frac{d}{dX}\Li_2(e^{-2X})=2\log(1-e^{-2X})$ and $\frac{d}{dX}\Li_2(-e^{-2X})=2\log(1+e^{-2X})$
(chain rule in $\Li_2'(z)=-\log(1-z)/z$), so the derivative of the right side is
$-\log\frac{1+e^{-2X}}{1-e^{-2X}}=\log\tanh X$; both sides vanish at $X=0$ because
$\Li_2(-1)=-\pi^2/12$, $\Li_2(1)=\pi^2/6$}
\lstep{1}{2}{$F_2(t)=8\int_0^{t/4}\log\tanh u\,du$ for $t>0$, hence $F_2'(t)=2\log\tanh(t/4)$ and
$F_2''(t)=\tfrac12\bigl(\sinh(t/4)\cosh(t/4)\bigr)^{-1}=1/\sinh(t/2)$.}
\lby{$\langle1\rangle1$ with $X=t/4$; $\frac{d}{dt}2\log\tanh(t/4)=\frac{2\cdot\frac14\sech^2(t/4)}{\tanh(t/4)}$}
\lstep{1}{3}{$\displaystyle\int_{\alpha_m}^{\beta_m}\!\!\int_{\alpha_n}^{\beta_n}\frac{d\xi\,d\xi'}{\sinh((\xi-\xi')/2)}
 =F_2(\beta_m-\alpha_n)-F_2(\alpha_m-\alpha_n)-F_2(\beta_m-\beta_n)+F_2(\alpha_m-\beta_n)$.}
\lby{$\langle1\rangle2$: with $F_1=F_2'$, $\int_{\alpha_n}^{\beta_n}f(\xi-\xi')d\xi'=F_1(\xi-\alpha_n)-F_1(\xi-\beta_n)$
and then integrate in $\xi$; the integral converges absolutely for disjoint $I_m,I_n$ (a logarithmic
singularity at a shared endpoint) and vanishes for $m=n$ by oddness}
\lstep{1}{4}{\eqref{eq:Qcell}.}
\lby{$\langle1\rangle3$, \eqref{eq:qt}, and $\int\!\!\int\chi_{I_m}\delta(\xi-\xi')\chi_{I_n}=h_m\delta_{mn}$}
\lqed{1}{5}
\end{proof}


\subsection{Validation: the discretised QFI form reproduces $f_2$}
\label{sec:conv}
Mesh: uniform cells of width $\approx h$ on $[\xi_0-\Lambda,\xi_1+\Lambda]$ with $\xi_0$ and $\xi_1$
exact grid points; $\hh_A,\hh_B,\hh_C$ are then exact coordinate subspaces of the Galerkin space
$\mathcal V_N=\mathrm{span}\{e_m\}$, which is what the box-Fourier frame of
\texttt{compression\_box.py} cannot provide.  $Q$ from Lemma~\ref{lem:closed}; the cell integrals of
$\widetilde\delta_c$ by tensor Gauss--Legendre ($10$ nodes), with a geometrically graded sub-rule
($24$ panels, ratio $0.55$) on the two cells touching the corner $\xi_0$, where the kernel
\eqref{eq:dc} is scale invariant.  Script \texttt{numerics/optimality\_all/kkt\_continuum.py},
output \texttt{kkt\_continuum\_conv.out}.  Weights are the \emph{true} $w^{\rm true}$; no cap is
needed, because the discrete $\kappa_{\max}$ is only $4.6$--$6.9$ (a piecewise-constant Galerkin
space is a poor resolver of the ultraviolet) and $w_{\max}\le\tfrac12(1+\cosh\kappa_{\max})\le500$.

\begin{table}[h]\centering\small
\begin{tabular}{rrrrrrrrr}
\toprule
$h$&$N$&$\kappa_{\max}$&$\dfrac{\|\delta_c\|_1}{\frac1{2\pi}\log(1+\zeta)}$&
$\dfrac{g_Q(\delta_c)}{f_2\zeta^2}$ (\texttt{first})&$\dfrac{g_Q(\delta_c)}{f_2\zeta^2}$ (\texttt{exact})\\
\midrule
0.400&  69&4.641&0.96673&0.86867&0.84819\\
0.300&  91&5.063&0.97701&0.89559&0.87565\\
0.250& 110&5.376&0.98426&0.91487&0.89553\\
0.200& 137&5.724&0.99016&0.93107&0.91240\\
0.150& 183&6.196&0.99674&0.94934&0.93171\\
0.125& 219&6.486&0.99980&0.95790&0.94090\\
0.100& 274&6.856&1.00315&0.96721&0.95105\\
\midrule
\multicolumn{4}{r}{Richardson $h\to0$ (\texttt{first}, last two / last three rows):}&
\multicolumn{2}{l}{$1.0045$ \quad / \quad $1.0030$}\\
\bottomrule
\end{tabular}
\caption{$s=0.1$, $\zeta=1/30$, $\Lambda=12$.  Convergence is clean first order,
$g_Q(\delta_c)/(f_2\zeta^2)=1-0.34\,h+O(h^2)$ over the whole range (the ratio
deviation$/h$ is $0.328,0.348,0.340,0.345,0.338,0.337,0.328$).  Infrared truncation is
irrelevant: at $h=0.25$, $\Lambda=12$ and $\Lambda=16$ give $0.91487$ and $0.91489$; at $h=0.2$,
$\Lambda=8$ gives $0.93022$ against $0.93107$ at $\Lambda=12$.  The \texttt{exact} column lies
$\approx2\%$ below the \texttt{first} column, the size expected from the cubic remainder
$c_3=-0.99$ at $\zeta=1/30$ combined with the $O(s^2)$ part of \eqref{eq:dc}.}
\label{tab:conv}
\end{table}

\begin{verdictok}[\;The modular cell discretisation is a faithful model of $g_Q(\delta_c)$]
Richardson extrapolation of Table~\ref{tab:conv} gives
$g_Q(\delta_c)/(f_2\zeta^2)=1.003\pm0.002$: the block-exact modular-frame Galerkin model
reproduces the continuum second-order value $f_2\zeta^2=\frac{c}{12\pi^2}\zeta^2$ of P2's
Theorem~A, independently of the box-Fourier computation of \texttt{compression\_box.py}.
Moreover the Legendre identity \eqref{eq:sldid} holds to $10^{-10}$ at every mesh
(\texttt{Legendre=1.0000000000} in \texttt{kc\_h05.out}, \texttt{kc\_h04.out}), so $t_c$, $V(t_c)$
and $g_Q(\delta_c)$ are internally consistent.
\end{verdictok}


\subsection{The linear programme at $t_c$, and the self-consistent KKT test}
\label{sec:LP}
Two SDPs are solved in the model of \S\ref{sec:disc}: the \emph{linear} programme
$\Delta(t_c)=\inf_{\mathcal F_N}\Tr t_c(Q-Q^{\rm rec})$, and the \emph{primal}
$\min_{\mathcal F_N}g_Q$ together with the linear programme at the SLD $t_{\rm opt}$ of its
minimiser.  Outputs \texttt{kc\_h05.out}, \texttt{kc\_h04.out}.

\begin{table}[h]\centering\small
\begin{tabular}{lcccccccc}
\toprule
mesh&$N$&$(n_A,n_B,n_C)$&$\dfrac{g_Q(\delta_c)}{f_2\zeta^2}$&$\dfrac{\Delta(t_c)}{2g_Q(\delta_c)}$&
$\dfrac{\Delta(t_c)_+^2/(8V)}{f_2\zeta^2}$&$\dfrac{\min_{\mathcal F_N}g_Q}{f_2\zeta^2}$&
$\dfrac{\min_{\mathcal F_N}g_Q}{g_Q(\delta_c)}$&$\dfrac{\Delta(t_{\rm opt})^2/(8V_{\rm opt})}{\min g_Q}$\\
\midrule
$h{=}0.5$, $\Lambda_A{=}8$, $\Lambda_C{=}4$&31&$(16,7,8)$&0.8359&$-9.395$&$\mathbf 0$&8.398&10.05&$1.0000073$\\
$h{=}0.4$, $\Lambda_A{=}8$, $\Lambda_C{=}6$&44&$(20,9,15)$&0.8675&$-10.325$&$\mathbf 0$&7.169&8.26&$1.0000186$\\
\bottomrule
\end{tabular}
\caption{$s=0.1$, $\zeta=1/30$, kernel \texttt{first}.  Column~5: the linear functional
$\Tr t_c(Q-Q^{\rm rec})$ is \emph{not} minimised at the compression --- its infimum over
$\mathcal F_N$ is about ten times $2g_Q(\delta_c)$ and of the opposite sign --- so the dual value it
certifies (column~6) is $0$, not $f_2\zeta^2$.  Column~9 is the same duality ratio computed
self-consistently at the model's \emph{own} optimum: it is $1$ to $2\times10^{-5}$, with KKT gaps
$-3.6\times10^{-6}$ and $-9.3\times10^{-6}$.}
\label{tab:LPtc}
\end{table}

\paragraph{Diagnosis: a finite Galerkin model cannot hold the compression.}
The zero-collar compression is the one-particle map of a Moebius \emph{squeeze} of $D=(0,L)$ onto
$B=(0,b)$; $X_c:\hh_B\to\hh_D$ is \emph{unitary}, so $X_cX_c^*=1_{\hh_D}$ and both constraints of
$\mathcal F$ are active.  In any finite model $\dim\mathcal V_B<\dim\mathcal V_D=\dim\mathcal V_B+\dim\mathcal V_C$,
hence for every $(X,Z)\in\mathcal F_N$
\begin{equation}\label{eq:deficit}
 \Tr(1_{\mathcal V_D}-XX^*)\ \ge\ \dim\mathcal V_C\ =\ n_C\ >\ 0 ,
\end{equation}
and $\delta_c\notin\{Q-Q^{\rm rec}(X,Z):(X,Z)\in\mathcal F_N\}$: the compression is \emph{infeasible}
in the truncated programme.  (Measured: $\Tr(1-XX^*)=11.66$ and $20.17$ at the optima of the two
rows of Table~\ref{tab:LPtc}, against $n_C=8,15$.)  This is fatal for the test at $t_c$ because
the KKT identity is a near-total cancellation: $\|t_c\|_1=1.3\times10^{-2}$ and the feasible set has
$O(1)$ diameter, so a generic $t$ gives $\Delta(t)=O(10^{-2})$, whereas
$2g_Q(\delta_c)=1.6\times10^{-5}$.  Since $|\Delta(t)-\Delta(t')|\le2\|t-t'\|_1$, a relative error of
only $10^{-3}$ in $t_c$ --- or in the feasibility of the point it is the SLD of --- already swamps
the answer.  In the modular coordinate the obstruction is quantitative: $\Xi:=\xi\circ h_s\circ\xi^{-1}$
maps $(\xi_0,\infty)$ onto $(\xi_0,\xi_1)$ with $\Xi(\xi)-\xi\simeq-\frac{s\mu}{L}(\xi-\xi_0)^2$
near the corner ($\mu=aL/(a+L)$), so the images $\Xi^k(C)$ accumulate at $\xi_0$ only like
$p_k\sim L/(s\mu k)$: representing the squeeze down to modular scale $p$ needs
$\sim L/(s\mu p)$ cells in $B$ against $O(1)$ cells for the physics.  \emph{The number of $B$-cells
needed grows like $1/s$ as the collar closes} --- precisely the regime the conjecture is about.

\paragraph{Consequence.}  The requested ratio $\Delta(t_c)^2/(8V(t_c))\big/(f_2\zeta^2)$ is
$\mathbf 0$ in every mesh tested, but this is a \emph{property of the truncation}, not evidence
against S10's KKT claim: the number is dominated by the infeasibility \eqref{eq:deficit} of the
compression, and moves by an order of magnitude when $n_B/n_C$ is changed at fixed physics
(Table~\ref{tab:s1}).  What the same computation \emph{does} establish, at machine precision, is
the duality theory itself: at the model's own optimum, $t_{\rm opt}$ is a certificate and
$\Delta(t_{\rm opt})^2/(8V(t_{\rm opt}))=\min_{\mathcal F_N}g_Q$ to $2\times10^{-5}$
(last column of Table~\ref{tab:LPtc}) --- Theorem~4.1(3) of \texttt{optimality\_all\_channels.tex}
verified in the continuum discretisation as well as on the lattice.


\subsection{How the answer depends on $\dim\mathcal V_B$ versus $\dim\mathcal V_C$}
\label{sec:s1}
To make the compression at least approximately representable one must refine $B$ far more than
$C$ (see the estimate above), and one must take the collar $s$ \emph{large}, since the number of
$B$-cells needed grows like $1/s$.  Table~\ref{tab:s1} therefore works at $s=1$, i.e.\ $b=L/2$,
$\zeta=1/3$, with region-dependent mesh widths $(h_A,h_B,h_C)$ and the \texttt{exact} kernel
\eqref{eq:dc}.  Script \texttt{scan\_s.py} (built on \texttt{kkt\_continuum.py}), outputs
\texttt{scan\_s1b.out}, \texttt{scan\_s1c.out}, \texttt{scan\_s1diag.out}.

\begin{table}[h]\centering\small
\begin{tabular}{cccccccccc}
\toprule
$h_A$&$h_B$&$h_C$&$N$&$(n_A,n_B,n_C)$&$\dfrac{g_Q(\delta_c)}{f_2\zeta^2}$&
$\dfrac{\min_{\mathcal F_N}g_Q}{f_2\zeta^2}$&$\dfrac{\min_{\mathcal F_N}g_Q}{g_Q(\delta_c)}$&
$\dfrac{\Delta(t_c)}{2g_Q(\delta_c)}$&$\dfrac{\Delta(t_{\rm opt})^2/(8V)}{\min g_Q}$\\
\midrule
0.50&0.35&1.500&24&$(16,\;4,\;4)$&0.6321&1.0245&1.6207&---&$1.000000$\\
0.50&0.24&1.500&26&$(16,\;6,\;4)$&0.6482&0.6507&1.0038&$-0.521$&$1.000001$\\
0.50&0.20&1.500&27&$(16,\;7,\;4)$&0.6535&0.5615&0.8591&---&$1.000002$\\
0.50&0.12&1.500&32&$(16,12,\;4)$&0.6684&0.4082&0.6108&$-0.874$&$1.000001$\\
0.50&0.12&0.750&36&$(16,12,\;8)$&0.7031&0.5666&0.8058&$-0.718$&$0.999998$\\
0.40&0.12&0.750&40&$(20,12,\;8)$&0.7192&0.5889&0.8188&$-0.647$&$0.999997$\\
\midrule
\multicolumn{5}{l}{\emph{for comparison, $s=0.1$ (Table~\ref{tab:LPtc}):}}&&&&&\\
0.50&0.50&0.500&31&$(16,\;7,\;8)$&0.8359&8.3980&10.047&$-9.395$&$1.0000073$\\
0.40&0.40&0.400&44&$(20,\;9,15)$&0.8675&7.1686&8.2635&$-10.325$&$1.0000186$\\
\bottomrule
\end{tabular}
\caption{$\Lambda_A=8$, $\Lambda_C=6$.  Reading the table: with $n_B<n_C$ the truncated programme
is \emph{far worse} than the compression ($\min/g(\delta_c)=8$--$10$), because \eqref{eq:deficit}
forces a rank-$n_C$ mismatch in the $A\!-\!D$ block; with $n_B\gg n_C$ it is \emph{better} than the
compression ($0.61$), because a coarse $C$ under-represents the hard part of the problem.  Along
the diagonal $n_B/n_C=1.5$ the ratio still drifts, $1.004\to0.819$.  The last column is the
self-consistent duality ratio at the model's own optimum: $1$ to $\le2\times10^{-5}$ in every row.}
\label{tab:s1}
\end{table}

\noindent Both truncation errors are one-sided and \emph{opposite}, and both are large compared
with the effect being tested.  Note also the two exact inequalities that make the model legitimate
but non-decisive: (i) $\mathcal F_N\subseteq\mathcal F$ --- every discrete feasible pair is a genuine
gauge-invariant quasi-free channel, by extending $X$ with $XP_{\mathcal V_B}$ and $Z$ by $0$ on
$\mathcal V_D^\perp$; and (ii) $g_{Q_N}(P\delta P)\le g_Q(\delta)$, because
$V_{Q_N}(t)=\Tr(tQt(1-Q))$ \emph{exactly} for $t$ supported on $\mathcal V_N$ (as
$P(1-Q)P=1_N-Q_N$), so $g_{Q_N}$ is the same Legendre supremum over a smaller set of test
operators.  The two inequalities push $\min_{\mathcal F_N}g_{Q_N}$ in opposite directions relative
to $\min_{\mathcal F}g_Q$, and neither can be dropped.


\subsection{The discrete minimiser is a truncation artefact (nested-refinement test)}
\label{sec:refine}
The rows of Table~\ref{tab:s1} with $\min_{\mathcal F_N}g_Q<g_Q(\delta_c)$ raise the question
whether a genuine channel beats the compression.  They do not, and this can be decided without a
larger SDP.  Split every cell of a mesh into $r$ equal sub-cells: then
$\mathcal V_N\subset\mathcal V_{N'}$, the inclusion $E:\mathcal V_N\to\mathcal V_{N'}$ is an isometry
mapping cells to normalised sums of sub-cells, and $(X,Z)\mapsto(E_DXE_B^*,\ E_DZE_D^*)$ maps
$\mathcal F_N$ into $\mathcal F_{N'}$ (the constraints are preserved because
$E_B^*Q_{N',BB}E_B=Q_{N,BB}$ and both slacks are $\ge0$ block by block).  Since
$g_{Q_{N'}}(P'\delta P')\uparrow g_Q(\delta)$, evaluating the coarse minimiser in the refinements
gives an increasing sequence of \emph{lower} bounds for its true continuum second-order error.
Script \texttt{refine\_check.py}, outputs \texttt{refine\_s1.out}, \texttt{refine\_s1b.out}.

\begin{table}[h]\centering\small
\begin{tabular}{rrcc|rrcc}
\toprule
\multicolumn{4}{c|}{coarse $(h_A,h_B,h_C)=(0.4,0.12,0.75)$, $(n_A,n_B,n_C)=(20,12,8)$}&
\multicolumn{4}{c}{coarse $(0.5,0.24,1.5)$, $(16,6,4)$}\\
$r$&$N'$&$\dfrac{g_{Q_{N'}}(\delta_c)}{f_2\zeta^2}$&
$\dfrac{g_{Q_{N'}}(\delta_{\rm opt})}{f_2\zeta^2}$&
$r$&$N'$&$\dfrac{g_{Q_{N'}}(\delta_c)}{f_2\zeta^2}$&
$\dfrac{g_{Q_{N'}}(\delta_{\rm opt})}{f_2\zeta^2}$\\
\midrule
1& 40&0.7192&0.5889          &1& 26&0.6482&0.6507\\
2& 80&0.7975&$6.40\e{3}$     &2& 52&0.7465&$3.11\e{3}$\\
3&120&0.8343&$2.01\e{4}$     &4&104&0.8193&$1.97\e{4}$\\
4&160&0.8578&$4.12\e{4}$     &8&208&0.8767&$9.58\e{4}$\\
6&240&0.8887&$1.06\e{5}$     & &   &      &\\
\bottomrule
\end{tabular}
\caption{$s=1$, $\zeta=1/3$.  The compression defect (columns 3,7) climbs monotonically towards
$f_2\zeta^2$, as it must.  The minimiser of the coarse programme (columns 4,8), embedded
\emph{unchanged} into the refinements, explodes by five orders of magnitude at the very first
refinement: as a genuine quasi-free channel its second-order error is
$\ge10^5\,f_2\zeta^2$.  It ``wins'' in the coarse model only because $\mathcal V_N$ cannot see the
modes into which it dumps its defect --- the uncapped continuum analogue of S8's ultraviolet
amplification (S8 Sec.~3.4).}
\label{tab:refine}
\end{table}

\begin{verdictgap}[\;(b) The certificate test is inconclusive, for a structural reason]
In the block-exact modular Galerkin model the requested ratio is
$\Delta(t_c)^2/(8V(t_c))\big/(f_2\zeta^2)=0$ in every mesh, because
$\Delta(t_c)<0$ (Tables~\ref{tab:LPtc},~\ref{tab:s1}).  This is \emph{not} evidence against S10's
KKT claim: the zero-collar compression has $X_cX_c^*=1_{\hh_D}$, which no finite-dimensional
Galerkin model can realise (eq.~\eqref{eq:deficit}), so $\delta_c$ is infeasible in
$\mathcal F_N$ and the near-total cancellation that KKT asserts (from
$\|t_c\|_1\sim10^{-2}$ down to $2g_Q(\delta_c)\sim10^{-5}$) cannot take place.  The measured
$\Delta(t_c)/(2g_Q(\delta_c))$ moves from $-10.3$ (at $n_B<n_C$) to $-0.52$ (at $n_B>n_C$),
i.e.\ it is governed by $\dim\mathcal V_B/\dim\mathcal V_C$ and not by the physics.  Symmetrically,
the readings $\min_{\mathcal F_N}g_Q<g_Q(\delta_c)$ are refuted by Table~\ref{tab:refine}: the
discrete minimisers are ultraviolet artefacts, not channels that beat the compression.  A
conclusive numerical test of the certificate needs a discretisation in which the Moebius squeeze
is exactly representable --- e.g.\ a mesh uniform in $\zeta=-L/x$, in which $h_s$ is the exact
translation $\zeta\mapsto\zeta-s$ --- at the price of a hard gap of size $L/\Zmax$ at the corner;
this is left as the next step.
\end{verdictgap}


\section{Part (c): exact root fidelity of S10's rotated channel}
\label{sec:s10}

\paragraph{The claim under test.}  Track S10 (\texttt{rigor/sdp\_dual\_certificate.tex}
Secs.~6,7,9) reports that in the spectral NS-circle model the zero-collar compression
$W_s=e^{-sD_\chi}$ is \emph{not} the second-order optimum: composing it with a Bogoliubov
rotation $e^{-\eps G_*}$ supported in $D$ lowers the SLD form by $1-\vartheta$,
$\vartheta=0.300\pm0.005$, the non-perturbative check giving the ratio $0.705$ at $\eps/s=-1$
(\texttt{s10\_nonpert.out}).  That evidence is at the level of the (capped) quadratic form
$g_Q$; an ``order of limits'' artefact --- $g_Q$ overestimating a defect that the exact fidelity
does not see --- is exactly what S8's Sec.~3.4 warns about.  We therefore compute the
\emph{exact} $-\log\Fid(\omega,\omega^\beta)$ in the same model.

\paragraph{Method.}  Script \texttt{numerics/optimality\_all/s11\_exactF.py} (outputs
\texttt{s11\_exactF\_L*.out}) rebuilds S10's ingredients through \texttt{s10\_theta.setup}
($P$ = spectral projection $\nu>0$ on the NS circle of $L$ sites; $D_\chi$ the tapered generator
of $\chi\partial_\theta$, $\chi=\cos\theta-1$ on $D$; $I=A\cup D$; $G_*$ from the same conjugate
gradient), and forms
\[
 Q=P|_I,\qquad Q^{\rm rec}(s,\eps)=\bigl(\widetilde W^*P\widetilde W\bigr)\big|_I,\qquad
 \widetilde W=e^{-sD_\chi}e^{-\eps G_*}.
\]
$-\log\Fid$ of the two quasi-free states with these symbols is Note~5 Thm.~2.1,
\begin{equation}\label{eq:note5}
 -\log\Fid=-\tfrac12\textstyle\sum_i\log(1-q_i)-\tfrac12\sum_j\log(1-q^{\rm rec}_j)
 -\sum_k\log\bigl(1+\sqrt{\tau_k}\bigr),\quad
 \tau=\spec\bigl(G_1^{1/2}G_2G_1^{1/2}\bigr),\ G=\tfrac{Q}{1-Q} .
\end{equation}
$Q=P|_I$ has eigenvalues exponentially close to $0$ and $1$, so \eqref{eq:note5} is a
catastrophic cancellation in double precision.  Following \texttt{numerics/fidelity\_flint.py} we
\emph{sub-compress} onto the spectral window
$\mathcal S(\kappa_c)=\mathrm{span}\{$eigenvectors of $Q$ with $|\kappa|=|\log\frac{1-q}{q}|\le\kappa_c\}$
and evaluate \eqref{eq:note5} there in \texttt{mpmath} at $60$ digits.  Compression is a channel,
so this is a rigorous \emph{lower} bound, increasing to the true value as $\kappa_c\uparrow$; both
channels are compressed with the \emph{same} $\mathcal S(\kappa_c)$, so the ratio compares like
with like.  Window sizes at $L=192$ ($|I|=114$): $\dim\mathcal S=12,16,20,24,26$ for
$\kappa_c=10,15,20,25,30$; beyond $\kappa_c\approx36$ the double-precision floor of $Q$ is reached.


\paragraph{Results.}
\begin{table}[h]\centering\small
\begin{tabular}{rlccccc|cc}
\toprule
&&\multicolumn{5}{c|}{$-\log\Fid$ ratio at $\eps/s=-1$, window $\kappa_c=$}&
\multicolumn{2}{c}{best over $\eps$}\\
$L$&$s$&$10$&$15$&$20$&$25$&$30$&$-\log\Fid$&$g_Q^{(\kappa_c=10^4)}$\\
\midrule
192&0.05&0.5889&0.6251&0.6490&0.6662&0.6720&$0.6692$ at $-1.10$&$0.7093$ at $-1.00$\\
192&0.02&0.5865&0.6207&0.6434&0.6597&0.6655&$0.6614$ at $-1.10$&$0.7049$ at $-1.00$\\
192&0.01&0.5889&0.6226&0.6450&0.6609&0.6666&$0.6624$ at $-1.10$&$0.7072$ at $-1.00$\\
320&0.02&---&0.6105&---&0.6488&$0.6660^{\ast}$&$0.6621$ at $-1.10$&$0.7052$ at $-1.00$\\
\bottomrule
\end{tabular}
\caption{Ratio $-\log\Fid(\omega,\omega^{\beta_{s,\eps}})\big/{-\log\Fid(\omega,\omega^{\beta_{s,0}})}$
of the rotated channel to the pure compression, evaluated by \eqref{eq:note5} on the spectral
window $\mathcal S(\kappa_c)$.  ${}^{\ast}$: $\kappa_c=32$ at $L=320$ ($\dim\mathcal S=31$;
the other two $L=320$ windows are $\kappa_c=15,25$ with $\dim\mathcal S=17,25$).  Absolute
denominators at $L=192$, $\kappa_c=30$:
$-\log\Fid=8.7380\e{-6},\ 1.4256\e{-6},\ 3.5852\e{-7}$ for $s=0.05,0.02,0.01$, against
$g_Q^{(10^4)}=9.0207\e{-6},\ 1.4499\e{-6},\ 3.6293\e{-7}$ (the windowed fidelity recovers
$97$--$99\%$ of the capped quadratic form, as it should at this $s$).  The last column reproduces
S10's non-perturbative number $0.705$ (\texttt{s10\_nonpert.out}) to $\pm0.005$, confirming that
the model has been rebuilt faithfully.  $\vartheta$ from the same conjugate gradient:
$0.28870$ at $L=192$, $0.29074$ at $L=320$.}
\label{tab:exactF}
\end{table}

\paragraph{Reading of Table~\ref{tab:exactF}.}
\begin{enumerate}[leftmargin=1.6em,itemsep=2pt]
\item \emph{The gain survives in the exact fidelity.}  In \emph{every} window, at every $s$ and
both $L$, the rotated channel has strictly smaller exact recovery error than the compression, by
at least $33\%$.  There is no sign of the ``order of limits'' artefact: the improvement is not an
artefact of the quadratic form.
\item \emph{Monotone in $\kappa_c$, and the limit is $0.67$--$0.71$.}  Since sub-compression is a
channel, each column is a rigorous lower bound for $-\log\Fid$ of \emph{both} states, and the
ratios rise monotonically with the window: $0.589\to0.625\to0.649\to0.666\to0.672$ at $s=0.05$.
The increments shrink ($0.036,0.024,0.017,0.006$), so the windowed ratio is close to saturation
near $0.67$; the remaining gap to the capped-quadratic value $0.705$ is of the size of the
modes above $\kappa_c\approx30$, which double-precision $Q$ cannot resolve.  A safe statement is
therefore \(\;-\log\Fid(\text{rotated})/{-\log\Fid}(\text{compression})=0.67\text{--}0.71\).
\item \emph{Stable in $s$ and in $L$.}  At $\kappa_c=30$ the ratio is $0.6720,0.6655,0.6666$ for
$s=0.05,0.02,0.01$: no drift as $s\to0$, so the number is the second-order ratio, not a
finite-$s$ effect.  At $s=0.02$, $L=192$ and $L=320$ give $0.6655$ and $0.6660$ in the same window
family ($\kappa_c=30$ vs $32$): no drift with the lattice cutoff either.
\item \emph{The exact optimum sits slightly past $\eps/s=-1$.}  Optimising $\eps$ on the exact
fidelity rather than on $g_Q$ moves the minimiser from $\eps/s=-1.00$ to $-1.10$ and improves the
ratio by a further $0.004$--$0.005$ --- a third-order effect, in the direction that makes the
compression \emph{worse}, not better.
\end{enumerate}

\begin{verdictbreak}[\;(c) S10's refutation is confirmed by the exact fidelity]
In S10's spectral NS-circle model the composition of the zero-collar compression with the
Bogoliubov rotation $e^{-\eps G_*}$, $\eps\simeq-s$, has exact recovery error
$-\log\Fid$ equal to $0.67$--$0.71$ times that of the compression, uniformly in $s\in[0.01,0.05]$
and in $L\in\{192,320\}$, computed from Note~5 Thm.~2.1 in $60$-digit arithmetic on spectral
windows $|\kappa|\le\kappa_c$ (each value a rigorous lower bound for the true $-\log\Fid$ of both
states).  The capped quadratic form of \texttt{s10\_nonpert.py} gives $0.705$ in the same runs.
Hence the $30\%$ gain is \emph{not} an artefact of the SLD form or of the weight cap: within this
model the zero-collar compression is not the second-order optimum among gauge-invariant
quasi-free channels, and therefore $\Erec^{(2)}\le(1-\vartheta)f_2z^2$ with $\vartheta\approx0.3$.
\end{verdictbreak}


\section{Verdict}
\label{sec:verdict}

\paragraph{(a) Lattice: strong duality and KKT hold, to solver precision.}
On the $(4,8,2)$ chain, for every cap $\kappa_c\in\{4,6,8,10,12,14\}$ at which Clarabel
converges, the SLD $t_{\rm opt}$ of the optimal defect of $\min_{\mathcal F}g_w$ minimises the
\emph{linear} functional $\Tr t_{\rm opt}(Q-Q^{\rm rec})$ over $\mathcal F$ (relative KKT gap
$3\e{-6}$ to $3\e{-4}$), and
$\Delta(t_{\rm opt})^2/(8V_w(t_{\rm opt}))=\min_{\mathcal F}g_w$ to $\le3\e{-4}$ relative
(Table~\ref{tab:lattice}).  This is Theorem~4.1(3) of \texttt{optimality\_all\_channels.tex}
(the Sion minimax step) verified numerically.  Two side findings: the Legendre pair must be kept
consistent --- using $V_{\rm true}=\Tr(tQt(1-Q))$ with a \emph{capped} $g_w$ produces a spurious
$1$--$7\%$ excess (Lemma~\ref{lem:legendre}) --- and S8's stored $(X,Z)$ in
\texttt{best\_4\_8\_2.npz} is $0.44\%$ above the true minimiser of its own capped form
($7.72397\e{-5}$ vs $7.69024\e{-5}$ at $\kappa_c=8$), which does not affect S8 Table~1 because
that table reports the exact $-\log\Fid$ evaluated at the point actually produced.

\paragraph{(b) Continuum: the certificate test at $t_c$ is inconclusive, for a structural reason.}
The block-exact modular Galerkin model of \S\ref{sec:disc} is validated
($g_Q(\delta_c)/(f_2\zeta^2)\to1.003\pm0.002$ by Richardson extrapolation, Table~\ref{tab:conv};
the Legendre identity holds to $10^{-10}$), and in it the duality theory is again confirmed at
machine precision at the model's own optimum (last columns of
Tables~\ref{tab:LPtc},~\ref{tab:s1}).  But the requested ratio is
\[
 \frac{\Delta(t_c)^2/(8V(t_c))}{f_2\zeta^2}\;=\;\mathbf 0\qquad\text{in every mesh tested,}
\]
because $\Delta(t_c)<0$: the compression's SLD is not a dual certificate \emph{of the truncated
programme}.  The reason is structural and is the main negative finding of this track: the
zero-collar compression is a Moebius \emph{squeeze} with $X_cX_c^*=1_{\hh_D}$, so in any finite
Galerkin model $\Tr(1-XX^*)\ge\dim\mathcal V_C>0$ (eq.~\eqref{eq:deficit}) and $\delta_c$ is
\emph{infeasible}; the KKT identity is a cancellation from $\|t_c\|_1\sim10^{-2}$ down to
$2g_Q(\delta_c)\sim10^{-5}$, which cannot survive that.  Empirically $\Delta(t_c)/(2g_Q(\delta_c))$
runs from $-10.3$ to $-0.52$ as $\dim\mathcal V_B/\dim\mathcal V_C$ is raised at fixed physics.
The apparent counter-readings $\min_{\mathcal F_N}g_Q<g_Q(\delta_c)$ are disposed of by the
nested-refinement test (Table~\ref{tab:refine}): embedded unchanged into a refinement, the coarse
minimisers have $g_Q\ge10^5f_2\zeta^2$ --- they are ultraviolet artefacts, not better channels.
A conclusive test needs a discretisation in which the squeeze is exact, e.g.\ uniform in
$\zeta=-L/x$ where $h_s$ is the translation $\zeta\mapsto\zeta-s$, at the cost of a hard corner
gap; that is the recommended next step, and it is \emph{not} what this track's numbers settle.

\paragraph{(c) The gain of S10's rotated channel is real.}
Because (b) leaves the certificate question open on the compression side, the decisive number is
the direct one: in S10's spectral NS-circle model the rotated channel
$e^{-sD_\chi}e^{-\eps G_*}$ at $\eps\simeq-s$ has \emph{exact} recovery error $0.67$--$0.71$ times
that of the compression, stable in $s\in[0.01,0.05]$ and in $L\in\{192,320\}$, each entry a
rigorous lower bound for both $-\log\Fid$'s (Table~\ref{tab:exactF}).  The capped quadratic form
gives $0.705$ on the same runs.  So S10's $\vartheta\approx0.3$ is not an artefact of the SLD
form or of the weight cap.  Combined with (a) --- which says that where the compression \emph{is}
the optimum its SLD must be the certificate --- the consistent picture is: in this model the
compression is \emph{not} the optimum, the certificate at $t_c$ therefore cannot exist, and the
negative $\Delta(t_c)$ of part (b), although quantitatively dominated by truncation, has the
right sign.  What remains unproved is whether the same holds in the continuum limit of the
physical geometry: (b)'s model cannot decide it, and (c)'s model carries the taper and the
lattice cutoff of \texttt{s10\_theta.setup}.

\paragraph{Deliverables.}
\texttt{numerics/optimality\_all/}: \texttt{kkt\_lattice.py}/\texttt{.out},
\texttt{kkt\_continuum.py} with \texttt{kkt\_continuum\_conv.out}, \texttt{kc\_h05.out},
\texttt{kc\_h04.out}; \texttt{scan\_s.py} with \texttt{scan\_s1b.out}, \texttt{scan\_s1c.out},
\texttt{scan\_s1diag.out}; \texttt{refine\_check.py} with \texttt{refine\_s1.out},
\texttt{refine\_s1b.out}; \texttt{s11\_exactF.py} with \texttt{s11\_exactF\_L192.out},
\texttt{s11\_exactF\_L320.out}; and this document.

\end{document}
